\documentclass[10pt,twocolumn]{article}
\usepackage[letterpaper,margin=0.75in]{geometry}
\usepackage{mathptmx}
\usepackage{booktabs}
\usepackage{amsmath,amssymb}
\usepackage{microtype}
\usepackage[colorlinks=true,allcolors=teal]{hyperref}
\setlength{\columnsep}{0.28in}
\setlength{\parskip}{2pt}
\renewcommand{\baselinestretch}{1.02}

\usepackage{titlesec}
\titleformat{\section}{\normalfont\large\bfseries}{\thesection}{0.6em}{}
\titleformat{\subsection}{\normalfont\normalsize\bfseries}{\thesubsection}{0.6em}{}
\titlespacing*{\section}{0pt}{8pt}{3pt}
\titlespacing*{\subsection}{0pt}{6pt}{2pt}

\newcommand{\Ab}{\bar{A}}

\begin{document}

\twocolumn[
\begin{center}
{\LARGE\bfseries Priority-Conditioned Retention: A Controllable\\[2pt]
Middle-Layer Memory Circuit in Pretrained Mamba\par}
\vspace{8pt}
{\large Yejun Ahn \quad Chanhee Jeong \quad Suyun Kim\par}
\vspace{3pt}
{\normalsize Department of Applied Artificial Intelligence, Sungkyunkwan University\\
2026-1 AI Systems --- Team Project (Proposal / Preliminary Report)\par}
\vspace{10pt}
\begin{minipage}{0.86\textwidth}
\small\noindent\textbf{Abstract.}
State-space models (SSMs) such as Mamba compress an entire history into a
fixed-size continuous hidden state, achieving linear-time inference at the cost
of associative recall: unlike attention, they cannot losslessly retrieve an
arbitrary past token. Prior work attributes this to a capacity bound but leaves
open whether the practical failure is a \emph{storage} limit or a \emph{control}
limit. We show that a \textbf{pretrained} Mamba-790m already contains a
\textbf{controllable, spatially localized memory circuit}, accessible at
inference with no fine-tuning. By scaling the pre-softplus input to the
$\Delta$ gate at positions downstream of a target, we reduce the effective decay
applied to that target and recover associative recall on a capacity-pressured
multi-query associative-recall (MQAR) task ($0.53 \to 0.73$ at the optimal
operating point). The effect is dose-dependent, target-specific, and
bidirectional, and it is \textbf{localized to layers 16--23}. Restricting the
intervention to these 8 of 48 layers recovers most of the recall gain at
\textbf{near-zero language-modeling cost (perplexity $\times 1.07$)}, whereas the
all-layer intervention inflates perplexity $\times 1.81$. We frame this as the
continuous-state analogue of importance-aware KV-cache compression and report an
honest boundary: a literal write-quantization analogue does \emph{not} transfer,
because the superimposed state has no per-item slots to reallocate.
\end{minipage}
\vspace{14pt}
\end{center}
]

\section{Introduction}
Mamba and related selective SSMs match Transformers on many language tasks while
running in linear time, but remain markedly weaker at \emph{associative recall}
--- retrieving a specific value bound to a key seen earlier in context. Jelassi
et al.~\cite{jelassi2024} show this stems from a fundamental capacity constraint:
a fixed state cannot losslessly copy arbitrarily long inputs. Yet a capacity
bound does not, by itself, explain \emph{behavioral} recall failure at moderate
load. The failure could be (a) a \textbf{control} problem --- the recurrence
applies a default recency decay and does not preferentially keep query-relevant
items --- or (b) a \textbf{storage} problem --- the state is genuinely out of
bits. These have opposite implications: (a) is fixable at inference; (b) is not.

This paper provides evidence for (a) in a regime where it matters. We do
\emph{not} modify or retrain the model. Instead we ask: does a pretrained Mamba
already expose a handle that lets us trade its default recency policy for a
target-preserving one? We find that it does, that the handle is the $\Delta$
gate, and --- most importantly --- that the handle is \textbf{localized to a
contiguous band of middle layers}. Operating only on that band recovers recall
essentially for free in terms of language-modeling quality, the behavior one
expects of a dedicated mechanism rather than a diffuse one.

\paragraph{Contributions.}
(1)~A test-time intervention on the $\Delta$ gate that controllably preserves a
target item in a capacity-pressured SSM, with dose-response, specificity, and
bidirectionality controls (\S\ref{sec:control}).
(2)~Localization of this control to \textbf{layers 16--23} of Mamba-790m, and a
\textbf{surgical} result: the 8-layer intervention matches most of the full-model
recall gain at perplexity $\times 1.07$ vs.\ $\times 1.81$
(\S\ref{sec:circuit}). This is our central finding.
(3)~A framing placing this control on a single axis with importance-aware
KV-cache compression, plus an honest negative result: a literal
write-quantization analogue fails, bounding how far the cache analogy carries
(\S\ref{sec:boundary}).
We report scope honestly: all positive results are for Mamba-790m on synthetic
MQAR; we observe non-replication on Mamba-1.4b and fragility on natural language
(\S\ref{sec:limits}).

\section{Background and Framing}
\paragraph{The $\Delta$ gate as a retention knob.}
In a selective SSM the state updates as
$h_t = \Ab_t h_{t-1} + \bar{B}_t x_t$, with $\Ab_t = \exp(\Delta_t A)$ and
$A<0$, so $\Delta_t>0 \Rightarrow \Ab_t \in (0,1)$. The contribution of token
$i$ surviving to a query at position $T$ is therefore
$\bar{B}_i x_i \cdot \prod_{\tau>i}\Ab_\tau$. Two facts follow. First, the
default policy is recency-biased: older contributions decay geometrically.
Second, the \emph{survival} of an item is governed by the \emph{downstream}
decays $\prod_{\tau>i}\Ab_\tau$, not by the item's own write alone. Reducing
$\Delta$ at positions after a target pushes those $\Ab_\tau \to 1$, slowing the
decay the target experiences --- a preservation lever requiring no weight change.

\paragraph{Discrete $\leftrightarrow$ continuous importance allocation.}
The KV-cache compression literature (KIVI~\cite{kivi}, KVQuant~\cite{kvquant},
ZipCache~\cite{zipcache}, and rotation-based quantizers such as
TurboQuant~\cite{turboquant}) implements one principle: spend more bits on more
important tokens. This presumes a slotted cache with random access. Mamba has no
slots --- all items are superimposed in one vector. We therefore treat
``importance $\to$ retained information'' as architecture-agnostic and
instantiate it in Mamba not as per-token bit allocation but as \emph{differential
retention} through the $\Delta$ gate. We name this framing Priority-Conditioned
Retention (PCR). \S\ref{sec:boundary} reports where the analogy holds and breaks.

\section{Method}
\paragraph{Task.}
Multi-query associative recall (MQAR) with random single-token key/value words
(no semantic priors). A prompt presents $N$ key--value pairs followed by a query
key; the model must emit the bound value. Capacity pressure is induced by
increasing $N$; we query the first pair as the default target.

\paragraph{Intervention.}
We register forward hooks on each layer's \texttt{dt\_proj} and multiply its
pre-softplus output by a per-position gain $g$. Setting $g<1$ (``other-gain'')
at positions downstream of the target's value reduces their effective decay. A
layer mask restricts the intervention to a chosen subset of layers. The model
runs on its pure-PyTorch sequential path (fused kernels intentionally not
installed) so the $\Delta$ pathway is hookable. No parameters are updated.

\paragraph{Model and metrics.}
\texttt{state-spaces/mamba-790m-hf} (48 layers, fp32). We report top-1 recall
over 30 random seeds per condition, and perplexity on a fixed held-out passage
under the same gain, relative to the unmodified model.

\section{Results}
\subsection{Test-time control}\label{sec:control}
Sweeping other-gain at $N{=}32$ over all layers yields a clear inverted-U with an
optimum at $g{=}0.8$: recall rises $0.53 \to 0.73$, then collapses for
$g \le 0.6$ as integration halts. The effect is \textbf{target-specific}:
applying $g{=}0.8$ to downstream positions yields $0.73$, while the same budget
on random positions yields only $0.50$ (near base), ruling out a generic
``more gain helps'' artifact.

\subsection{Layer localization and the surgical circuit}\label{sec:circuit}
Restricting $g{=}0.8$ to 8-layer bands isolates the responsible depth: recovery
concentrates in \textbf{L16--23} ($\Delta\,{+}0.13$), with a secondary
contribution at L24--31 ($+0.07$) and nothing elsewhere. Crucially, the
localization is also \emph{cheaper}: Table~\ref{tab:main} shows the 8-layer
intervention captures most of the recall gain at near-zero language-modeling
cost, while the all-layer version pays a large perplexity tax. A targeted edit to
the memory circuit is almost free; a global one is not.

\begin{table}[t]
\centering\small
\caption{Test-time intervention on Mamba-790m, MQAR $N{=}32$
(base recall $0.53$, base ppl $\approx 232$).}
\label{tab:main}
\begin{tabular}{lcc}
\toprule
Intervention & Recall & Perplexity \\
\midrule
All layers, $g{=}0.8$ & 0.73 & $\times 1.81$ \\
\textbf{L16--23 only}, $g{=}0.8$ & \textbf{0.67} & $\mathbf{\times 1.07}$ \\
\bottomrule
\end{tabular}
\end{table}

\section{The Discrete--Continuous Boundary}\label{sec:boundary}
If retention is one lever, \emph{precision} --- the literal KV-quant analogue ---
is a candidate second lever: quantize the write of low-priority tokens to fewer
bits. We tested this by quantizing the per-position convolution output (the
token's written representation) to $k$ bits. Single-item precision is not a
graceful knob: target recall is flat from $16\to 8$ bits, then cliffs
($8$-bit $0.53 \to 4$-bit $0.03 \to 0$). Reallocation fails outright: quantizing
\emph{distractor} positions to 4 bits drives target recall to $0.00$ --- it
corrupts the shared computation rather than freeing capacity.

This is informative, not merely a null. Because the state is a superposition with
no random-access slots, lowering one item's write precision either does nothing
($\ge 8$-bit) or injects noise the whole sequence shares ($\le 4$-bit). There is
no per-item bit budget to reallocate. The importance-allocation principle thus
transfers to Mamba \emph{only} through retention (\S\ref{sec:circuit}), not
through quantization --- a concrete boundary on the KV-cache analogy and an
empirical instance of the ``no slots'' property that distinguishes
continuous-state memory from a discrete cache.

\section{Limitations and Scope}\label{sec:limits}
We state these plainly; they bound the claim and we believe reporting them
strengthens it. \textbf{Single model:} all positive results are for Mamba-790m;
on Mamba-1.4b the same intervention shows no positive circuit at a comparable
pressure regime (all layer-band $\Delta \le 0$ despite ample headroom), so the
L16--23 circuit appears model-specific and we do not claim a universal depth.
\textbf{Synthetic task:} effects are robust on random-token MQAR but fragile on
natural-language multi-fact recall, where the intervention can hurt.
\textbf{Graded priority is not established:} a three-level graded variant
appeared monotone but is confounded with query position; we do not claim
multi-level allocation. The SSM constraint
($\Delta>0 \Rightarrow \Ab<1$, shared decay) predicts that truly item-selective
allocation is limited, consistent with \S\ref{sec:boundary}.

\section{Related Work}
\textbf{SSM capacity.} Jelassi et al.~\cite{jelassi2024} establish the
copying/capacity limit; work on primacy/recency documents the default decay
policy. We add a controllability result on top of the capacity picture.
\textbf{KV-cache compression.} KIVI, KVQuant, ZipCache and rotation-based
quantizers (TurboQuant~\cite{turboquant}) realize importance-aware bit allocation
on slotted caches; we ask whether the principle survives the move to a slotless
continuous state, and find it survives as retention but not as quantization.
\textbf{Recall and the ``what to keep'' frontier.} Prompting methods such as
Just Read Twice~\cite{jrt} and architectures such as Based, DeltaNet and
Titans~\cite{titans} target the same problem, mostly by training new
architectures or memory modules; ours is a training-free, frozen-base
alternative. \textbf{Mechanistic interpretability.} Localizing a behavior to a
contiguous layer band parallels circuit-style findings in Transformers; ours is,
to our knowledge, an early such localization of a \emph{controllable} memory
mechanism in a pretrained SSM.

\section{Conclusion and Future Work}
A pretrained Mamba's associative recall is, in part, a controllable property of a
specific middle-layer band, not a fixed ceiling. A surgical, training-free
intervention on layers 16--23 recovers recall at negligible language-modeling
cost. The importance-aware compression principle transfers to continuous state
through retention but not through quantization. Our proposed next step is to
remove the oracle: a small \emph{question-conditioned, auto-localizing}
controller trained on a frozen base, evaluated against Belady-MIN retention
efficiency under streaming queries, toward a training-free improvement of the
recall--capacity frontier of pretrained SSMs.

\begin{thebibliography}{9}\small
\bibitem{jelassi2024} S.~Jelassi et al. Repeat After Me: Transformers are Better than SSMs at Copying. \emph{ICML}, 2024.
\bibitem{kivi} Z.~Liu et al. KIVI: A Tuning-Free Asymmetric 2bit Quantization for KV Cache. \emph{ICML}, 2024.
\bibitem{kvquant} C.~Hooper et al. KVQuant: Towards 10M Context Length LLM Inference. \emph{NeurIPS}, 2024.
\bibitem{zipcache} Y.~He et al. ZipCache: Accurate and Efficient KV Cache Quantization with Salient Token Identification. \emph{NeurIPS}, 2024.
\bibitem{turboquant} A.~Zandieh et al. TurboQuant: Online Vector Quantization with Near-Optimal Distortion Rate. \emph{ICLR}, 2026.
\bibitem{jrt} S.~Arora et al. Just Read Twice: Closing the Recall Gap for Recurrent Language Models. arXiv:2407.05483, 2024.
\bibitem{titans} A.~Behrouz et al. Titans: Learning to Memorize at Test Time. 2024.
\bibitem{mamba} A.~Gu and T.~Dao. Mamba: Linear-Time Sequence Modeling with Selective State Spaces. arXiv:2312.00752, 2023.
\end{thebibliography}

\end{document}
